\section{Discussion}

Results show a clear trade-off between predictive accuracy and deployment constraints. Although Random Forest reached slightly lower MAE, the SNN remained competitive while offering event-driven temporal processing and direct alignment with neuromorphic edge deployment goals. This is relevant for low-power autonomous monitoring where model choice is not driven by accuracy alone. From a methodological standpoint, the multimodal setup proved valuable because it keeps physically related targets in a single inference pipeline and simplifies downstream decision logic in field applications. The observed error profile also indicates nonuniform target difficulty: pH and OD were more stable, while turbidity concentrated a larger share of uncertainty. This asymmetry is expected in heterogeneous freshwater environments and supports the use of weighted losses, repeated-seed reporting, and per-target diagnostics rather than relying on a single aggregate metric.

From the systems perspective, platform behavior is heterogeneous: Device B reached the highest throughput, Device C achieved the lowest incremental energy per inference, and ESP32 provided full offline feasibility at much lower processing rates. Therefore, deployment decisions should be scenario-dependent: high-rate analytics benefit from smartphone/SBC-class hardware, while sparse periodic sensing can still be addressed by microcontrollers. Energy interpretation should be made with explicit operational context, since baseline platform power varies with architecture and operating-system behavior. For this reason, energy per inference combined with throughput is a more actionable criterion than absolute watts alone. The environmental interpretation is consistent with the literature and regulatory practice, but generalization remains limited by dataset size. Augmentation and multi-seed evaluation reduced split sensitivity, yet broader field campaigns are still required for stronger external validity. Offline inference remains a practical advantage in low-connectivity regions by preserving local prediction and screening without cloud dependence.
